\documentclass[12pt]{article}
\usepackage[T1]{fontenc}
\usepackage[utf8]{inputenc}
\usepackage[brazil]{babel}
\usepackage[margin=1in]{geometry}
\usepackage{ragged2e}
\usepackage[normalem]{ulem}
\setlength{\parindent}{0pt}
\setlength{\parskip}{0.8\baselineskip}
\begin{document}
\justifying
\noindent Verifica-se que a questão objeto da controvérsia jurídica suscitada no recurso manejado pela parte recorrente esteve afetada no representativo da controvérsia - \textbf{Tema n. }\textbf{273} da TNU - (PU no processo n. PEDILEF 0043092-25.2017.4.03.6301/SP),  afetado em 21/08/2020 foi julgado (trânsito em julgado em Sobrestado em Secretaria - Tema 1220/STJ), por meio do qual foi elucidada a seguinte questão:

\noindent \textit{"Se particular que move ação própria, em decorrência de reconhecimento administrativo, operado em ação coletiva, da qual não fez parte, está jungido aos termos do acordo lá realizado."}

\noindent Restou firmada a seguinte tese:

\noindent \textit{"(i) no que toca à revisão do art. 29, II, da Lei 8.213/91, não é possível, valendo-se do título judicial formado na ação civil pública nº 0002320-59.2012.4.03.6183, inclusive dos valores em decorrência dele apurados, intentar ação para cumprimento do julgado (execução) com o objetivo de pagamento imediato, sem observância do cronograma estabelecido; (ii) o beneficiário do RGPS pode mover ação individual para revisão e/ou pagamento de parcelas vencidas decorrentes da correta aplicação do art. 29, II, da Lei 8.213/91, sem qualquer vinculação restritiva ao decidido na ação civil pública nº 0002320-59.2012.4.03.6183, inclusive no que toca ao cronograma de pagamento; (iii) intentada a ação individual, a contagem dos prazos de decadência do direito de revisão e da prescrição das parcelas vencidas deve observar o disposto no tema 134 da TNU." (Tema 273 TNU).}

\noindent Com efeito, depreende-se da leitura do voto condutor do acórdão que o mesmo se encontra em conformidade com o decidido no  tema pela TNU.

\noindent A proposito, consta do acordao hostilizado: (...) Anuidades OAB, Conselhos Regionais de Fiscalização Profissional e Afins, Organização Político-administrativa / Administração Pública, DIREITO ADMINISTRATIVO E OUTRAS MATÉRIAS DE DIREITO PÚBLICO, Extinção da Execução, Liquidação / Cumprimento / Execução, DIREITO PROCESSUAL CIVIL E DO TRABALHO, Nulidade / Inexigibilidade do Título, Liquidação / Cumprimento / Execução, DIREITO PROCESSUAL CIVIL E DO TRABALHO

\noindent Nessas condições, o conteúdo do Acórdão recorrido é consentâneo com o entendimento da TNU, firmado em julgamento de recurso representativo de controvérsia, o que atrai a incidência da regra prevista pelo art. 14, III, \textit{b,} da Resolução CJF n. 586/2019:

\noindent \textit{\textbf{Art. 14.}}\textit{ Decorrido o prazo para contrarrazões, os autos serão conclusos ao magistrado responsável pelo exame preliminar de admissibilidade, que deverá, de forma sucessiva: (...) }\textit{\textbf{III }}\textit{- negar seguimento a pedido de uniformização de interpretação de lei federal interposto contra acórdão que esteja em conformidade com entendimento consolidado: }\textit{\textbf{b)}}\textit{ em }\uline{\textit{\textbf{recurso representativo de controvérsia pela Turma Nacional de Uniformização}}}\textit{ ou em pedido de uniformização de interpretação de lei dirigido ao Superior Tribunal de Justiça;(...).}

\noindent Posto isso, com arrimo no \uline{\textbf{art. 14, III, }}\uline{\textit{\textbf{b}}}\textit{,} da Resolução CJF n. 586/2019 c/c art. 1.030, I, do CPC, \textbf{NEGO SEGUIMENTO }\textbf{a}o\textbf{ PU}\textbf{.}

\noindent Intimem-se. Aguarde-se o decurso do prazo recursal. Após, certifique-se o trânsito em julgado e devolva-se ao juízo de origem.
\end{document}
